Una corredera de un banco de ensayo idealizado se mueve horizontalmente y está unida a un resorte de rigidez \(k=4\,\mathrm{N/m}\). Su masa es \(m=1\,\mathrm{kg}\). El desplazamiento \(x\), en metros, se mide hacia la derecha desde la posición en la que el resorte no está deformado. El resorte ejerce la fuerza \(-kx\).

Se adopta esta ley de rozamiento: mientras la corredera desliza, la fuerza de rozamiento tiene módulo \(1\,\mathrm N\) y sentido contrario a la velocidad; en reposo puede tomar cualquier valor de módulo no superior a \(1\,\mathrm N\). Si esa fuerza puede equilibrar la del resorte, la corredera permanece en reposo; si no puede, empieza a deslizar en el sentido de la fuerza elástica. No actúan otras fuerzas horizontales.

En \(t=0\) se retira, sin impulso, la sujeción que mantenía la corredera en \(x=0.60\,\mathrm m\), en reposo. Los datos y las leyes son sintéticos. Formule el modelo de cada fase que llegue a producirse y obtenga la evolución para \(t\ge0\). Localice la primera parada posterior a la liberación y decida qué ocurre después. Compruebe que la variación de energía mecánica concuerda con el trabajo del rozamiento.
